\documentclass[12pt,a4paper]{article}
\usepackage[margin=24mm]{geometry}
\usepackage{fontspec}
\defaultfontfeatures{Ligatures=TeX}
\setmainfont{Liberation Serif}
\setsansfont{Liberation Sans}
\usepackage{polyglossia}
\setdefaultlanguage{russian}
\setotherlanguage{english}
\newfontfamily\cyrillicfont{Liberation Serif}
\usepackage{amsmath,amssymb,amsthm}
\usepackage{microtype}
\usepackage[unicode,hidelinks]{hyperref}
\newtheorem{theorem}{Теорема}[section]
\newtheorem{lemma}{Лемма}[section]
\newtheorem{definition}{Определение}[section]
\setlength{\emergencystretch}{3em}
\setlength{\parskip}{3pt}
\linespread{1.05}
\begin{document}
\begin{center}
{\LARGE\bfseries Исправленные формулировки и доказательства\par}
\medskip
{\large Комбинаторика слов и символическая динамика\par}
\medskip
17 ноября 2026 года
\end{center}
Этот текст содержит законченные замены красных, оранжевых и янтарных
узлов экспертной карты. Они также внесены в полную исправленную
редакцию диссертации. Номера утверждений в этом отдельном документе
локальны; соответствие исходной рукописи указано в заголовках разделов.
Исправление определения семейства не единственно: выбрана нормировка,
сохраняющая при $p=1$ исходную производящую функцию и порядок $k$.
Классическая классификация штурмовых степеней явно указана как
внешний результат; вывод асимптотического показателя из неё дан ниже.
\section{Равномерное распределение: исходная лемма 2.2.1}
\begin{theorem}[Критерий Вейля]
Последовательность $(x_n)$ равномерно распределена по модулю $1$
тогда и только тогда, когда для каждого целого $h\ne0$
$N^{-1}\sum_{n=1}^N\exp(2\pi\mathrm i h x_n)\to0$.
\end{theorem}
Последовательность логарифмов факториалов обладает более сильным
свойством, чем плотность множества её дробных частей: она равномерно
распределена по модулю $1$. Для десятичного основания это классический
результат о ведущих цифрах факториалов~\cite{Diaconis1977};
ниже приведено доказательство для любого фиксированного основания.

\begin{lemma}\label{lem:logfact}
Пусть $b>1$. Последовательность $(\log_b(n!))_{n\ge1}$
равномерно распределена по модулю $1$. В частности, для каждого
$[\alpha,\beta)\subseteq[0,1)$ при $\alpha<\beta$ существует
бесконечно много $n$ таких, что
$\{\log_b(n!)\}\in[\alpha,\beta)$.
\end{lemma}
\begin{proof}
Зафиксируем целое $h\ne0$ и положим
\[
 u_n=\exp(2\pi\mathrm i h\log_b(n!)),\qquad
 \tau=\frac{2\pi h}{\log b}\ne0.
\]
Для каждого фиксированного целого $r\ge1$ имеем
\[
 u_{n+r}\overline{u_n}
 =\exp\left(\mathrm i\tau\sum_{j=1}^r\log(n+j)\right)
 =n^{\mathrm i\tau r}+O_{b,h,r}(n^{-1}).
\]
Здесь использовано равенство
$\log(n+j)=\log n+O_r(n^{-1})$.
Для фиксированного вещественного $a$ сравнение суммы с интегралом даёт
\[
 \sum_{n=1}^M n^{\mathrm i a}
 =\frac{M^{1+\mathrm i a}}{1+\mathrm i a}+O_a(1+\log M).
\]
Действительно, ошибка сравнения ограничена величиной
$1+\int_1^M |(x^{\mathrm i a})'|\,dx
\le1+|a|\log M$, а интеграл вычисляется непосредственно.
Следовательно, для
$C_N(r)=\sum_{n=1}^{N-r}u_{n+r}\overline{u_n}$ получаем
\[
 \limsup_{N\to\infty}\frac{|C_N(r)|}{N}
 =\frac1{\sqrt{1+\tau^2r^2}}.
\]

Приведём необходимую оценку средних, чтобы явно указать все шаги.
Продолжим $u_n$ нулями за пределы $1\le n\le N$.
Для фиксированного целого $H\ge2$ имеем
\[
 H\sum_{n=1}^N u_n
 =\sum_{m=1}^{N+H-1}\sum_{j=0}^{H-1}u_{m-j}.
\]
Неравенство Коши--Буняковского и раскрытие квадрата дают
\[
 H^2\left|\sum_{n=1}^N u_n\right|^2
 \le (N+H-1)\left(
 HN+2\sum_{r=1}^{H-1}(H-r)\operatorname{Re}C_N(r)\right).
\]
Разделив на $H^2N^2$ и сначала переходя к верхнему пределу по $N$,
получаем
\[
\begin{aligned}
 \limsup_{N\to\infty}
 \left|\frac1N\sum_{n=1}^N u_n\right|^2
 &\le \frac1H+
 \frac2H\sum_{r=1}^{H-1}
 \left(1-\frac rH\right)\frac1{\sqrt{1+\tau^2r^2}}\\
 &\le \frac1H+\frac{2}{|\tau|H}\sum_{r=1}^{H-1}\frac1r
 \xrightarrow[H\to\infty]{}0.
\end{aligned}
\]
По критерию Вейля последовательность равномерно распределена.
Доля попаданий в интервал положительной длины стремится к
$\beta-\alpha>0$, поэтому число таких попаданий бесконечно.
\end{proof}
\section{Согласованная рекурсия: исходное определение 3.1.2}
Семейства обобщённых последовательностей Фибоначчи, зависящие
от параметров, рассматриваются, например, в~\cite{Mehraban2006Gulliver}.
Чтобы однозначно зафиксировать индексацию, в настоящей работе
используется следующая нормировка. Она выбрана так, чтобы при $p=1$
порядок рекурсии был равен $k$.

\begin{definition}\label{defn1pil}
Пусть $k\ge2$, $p\ge1$, $t\ge1$ --- целые числа и $r=k+p-1$.
Обобщённая $(t,k)$-$p$-последовательность
$(F_n^p(t,k))_{n\ge0}$ определяется равенствами
\begin{equation}
 F_n^p(t,k)=tF_{n-1}^p(t,k)
       +\sum_{j=p+1}^{r}F_{n-j}^p(t,k),
 \qquad n\ge r,
\end{equation}
и начальными условиями
\begin{equation}
 F_0^p(t,k)=\cdots=F_{r-2}^p(t,k)=0,
 \qquad F_{r-1}^p(t,k)=1.
\end{equation}
При необходимости полагаем $F_n^p(t,k)=0$ для $n<0$.
\end{definition}

Начальные условия задают ровно $r$ первых членов, а рекурсия
определяет каждый следующий член через ранее заданные.
По индукции последовательность существует и единственна.
В частности, при $p=1$ начальные условия имеют вид
$F_0^1=\cdots=F_{k-2}^1=0$, $F_{k-1}^1=1$, и
\[
 F_n^1(t,k)=tF_{n-1}^1(t,k)+F_{n-2}^1(t,k)+\cdots+F_{n-k}^1(t,k)
 \qquad(n\ge k).
\]
При $t=1$, $k=2$, $p=1$ получается обычная последовательность
Фибоначчи с $F_0=0$, $F_1=1$.
\section{Производящая функция: исходная лемма 3.1.1}
\begin{lemma}\label{lemgenn1}
Для последовательности из определения~\ref{defn1pil}, при $r=k+p-1$,
формальная производящая функция равна
\[
 \sum_{n\ge0}F_n^p(t,k)x^n
 =\frac{x^{r-1}}{1-tx-\sum_{j=p+1}^{r}x^j}.
\]
В частности, при $p=1$
\begin{equation}
 g_{F_n^1(t,k)}(x)
 =\sum_{n\ge0}F_n^1(t,k)x^n
 =\frac{x^{k-1}}{1-tx-x^2-\cdots-x^k}.
\end{equation}
\end{lemma}
\begin{proof}
Положим $G(x)=\sum_{n\ge0}F_n^p(t,k)x^n$ и
$D(x)=1-tx-\sum_{j=p+1}^{r}x^j$.
Соглашение о нулевых членах с отрицательными индексами даёт
\[
 [x^n](D(x)G(x))
 =F_n^p(t,k)-tF_{n-1}^p(t,k)
       -\sum_{j=p+1}^{r}F_{n-j}^p(t,k).
\]
При $n<r-1$ этот коэффициент равен $0$ по начальным условиям;
при $n=r-1$ он равен $1$; при $n\ge r$ он равен $0$
по рекурсии. Поэтому $D(x)G(x)=x^{r-1}$.
Свободный член $D$ равен $1$, следовательно, $D$ обратим
в кольце формальных степенных рядов. Это доказывает общую формулу.
При $p=1$ имеем $r=k$, что даёт указанное частное выражение.
\end{proof}
\section{Индексы штурмовых слов: исходная теорема 3.1.1}
Зафиксируем используемые понятия индекса. Если $v$ --- непустое
конечное слово и $s\ge1$ --- рациональное число с $s|v|\in\mathbb N$,
то $v^s$ означает префикс слова $v^\omega$ длины $s|v|$.
Для бесконечного слова $u$ положим
\[
 \operatorname{ind}_u(v)
 =\sup\{s\in\mathbb Q:\ s\ge1,\ s|v|\in\mathbb N,\ v^s\in\mathcal L_u\},
 \]
 \[
 E(u)=\sup_{v\in\mathcal L_u\setminus\{\varepsilon\}}\operatorname{ind}_u(v).
\]
Асимптотический критический показатель определяется как
\[
 E^*(u)=\lim_{M\to\infty}
 \sup_{\substack{v\in\mathcal L_u\\ |v|\ge M}}\operatorname{ind}_u(v).
\]
Предел существует в $[1,+\infty]$, поскольку выражение под
пределом не возрастает по $M$. В дальнейшем запись
$\operatorname{ind}(u)$ означает $E(u)$, а
$\operatorname{ind}^*(u)$ --- именно $E^*(u)$.

Перестановка букв сохраняет показатели $E$ и $E^*$.
Поэтому для штурмова слова с угловым коэффициентом $\alpha$
можно нормировать алфавит так, чтобы используемый коэффициент
был $\beta=\min\{\alpha,1-\alpha\}<1/2$.
В следующем утверждении $\varphi=(1+\sqrt5)/2$.

\begin{theorem}[Классические формулы для штурмовых слов
{\cite{DamanikLenz2002,Peltomaki2015,AdamczewskiAllouche2007}}]\label{thmconsn1}
Пусть $u$ --- штурмово слово с угловым коэффициентом
$\alpha\in(0,1)\setminus\mathbb Q$.
Положим $\beta=\min\{\alpha,1-\alpha\}=[0;a_1,a_2,\ldots]$,
$q_{-1}=0$, $q_0=1$ и
$q_{n+1}=a_{n+1}q_n+q_{n-1}$.
Тогда
\begin{align}
 E(u)&=\max\left\{a_1,\
       \sup_{n\ge1}\left(2+a_{n+1}+\frac{q_{n-1}-2}{q_n}\right)\right\},
       \label{eq:correct-index}\\
 E^*(u)&=2+\limsup_{n\to\infty}\frac{q_{n+1}}{q_n}
       =2+\limsup_{n\to\infty}[a_n;a_{n-1},\ldots,a_1].
       \label{eq:correct-asymptotic-index}
\end{align}
В частности, $E^*(u)\ge2+\varphi$.
\end{theorem}
\begin{proof}
Поскольку обмен букв не меняет показатели, работаем с нормированным
языком углового коэффициента $\beta<1/2$.
Используем следующую классическую структурную информацию о степенях
в штурмовом языке
\cite[теорема~4.5, следствие~4.6 и доказательство теоремы~4.7]{Peltomaki2015}:
\begin{itemize}
 \item если примитивное слово $v$ (не являющееся нетривиальной
 степенью более короткого слова) удовлетворяет
 $v^3\in\mathcal L_u$, то либо $v=0$, либо $|v|=q_n$
 для некоторого $n\ge1$;
 \item $\operatorname{ind}_u(0)=a_1$, а наибольший дробный показатель
 среди примитивных факторов длины $q_n$ равен
 \[
 e_n=2+a_{n+1}+\frac{q_{n-1}-2}{q_n}
     =2+\frac{q_{n+1}}{q_n}-\frac2{q_n}.
 \]
 Для каждого $n\ge1$ соответствующий фактор и его степень
 действительно встречаются.
\end{itemize}
Это используемый внешний комбинаторный результат; последующий вывод
асимптотической формулы приводится полностью.

Сначала заметим, что
$L=\limsup_{n\to\infty}q_{n+1}/q_n\ge\varphi$.
Если бы $L<\varphi$, можно было бы выбрать $c$ с
$L<c<\varphi$. Для всех достаточно больших $n$
$q_{n+1}/q_n\le c$, но рекурсия тогда даёт
\[
 \frac{q_{n+2}}{q_{n+1}}
 =a_{n+2}+\frac{q_n}{q_{n+1}}
 \ge1+\frac1c>c,
\]
что противоречит выбору $c$. Поэтому, поскольку $q_n\to\infty$,
\[
 \limsup_{n\to\infty}e_n=2+L\ge2+\varphi>3.
\]

Каждое непустое слово $w$ единственным образом записывается
$w=v^t$, где $v$ примитивно, а $t\ge1$ --- целое.
Из определения степеней следует
$\operatorname{ind}_u(w)\le\operatorname{ind}_u(v)/t$.
Если $\operatorname{ind}_u(w)>3$, то $v^3$ встречается.
По структурному результату $v=0$ или $|v|=q_n$;
показатель не превосходит соответственно $a_1$ или $e_n$.
Показатели $\le3$ не могут увеличить супремум, поскольку
$\sup e_n>3$. Обратные оценки обеспечиваются существованием
$0^{a_1}$ и степеней с показателями $e_n$.
Это доказывает~\eqref{eq:correct-index}.

Для асимптотической формулы существование степеней с корнями
длины $q_n\to\infty$ даёт
$E^*(u)\ge\limsup e_n=2+L$.
Если $L=+\infty$, равенство уже доказано.
В противном случае рассмотрим любую последовательность факторов
$w_j$ с $|w_j|\to\infty$ и показателями, стремящимися к $E^*(u)$.
Достаточно рассматривать показатели $>3$.
Для примитивного корня $v_j$ слова $w_j$ тогда имеем
либо $v_j=0$, либо $|v_j|=q_{n_j}$.
Первый случай не может встречаться при неограниченных $|w_j|$:
длина постоянного фактора не превосходит $a_1$.
Если индексы $n_j$ ограничены, то для каждого такого индекса
показатель примитивного корня конечен, и условие
$\operatorname{ind}_u(w_j)>3$ ограничивает также
$t_j=|w_j|/|v_j|$ величиной $e_{n_j}/3$.
Это вновь противоречит $|w_j|\to\infty$.
Следовательно, $n_j\to\infty$ после выбора подпоследовательности, и
\[
 \limsup_j\operatorname{ind}_u(w_j)
 \le\limsup_j e_{n_j}\le2+L.
\]
Таким образом, $E^*(u)=2+L$.
Наконец, рекурсия знаменателей даёт зеркальную формулу
$q_{n+1}/q_n=[a_{n+1};a_n,\ldots,a_1]$;
сдвиг индекса верхнего предела доказывает последнее равенство
в~\eqref{eq:correct-asymptotic-index}.
\end{proof}

\noindent Для слова Фибоначчи
$\beta=\varphi^{-2}=[0;2,1,1,\ldots]$ и $q_n=F_{n+2}$.
Из формулы Бине
$F_{n+3}-\varphi F_{n+2}=\psi^{n+2}$, где $\psi=-\varphi^{-1}$,
следует
\[
 e_n=2+\varphi+\frac{\psi^{n+2}-2}{F_{n+2}}<2+\varphi,
 \qquad e_n\longrightarrow2+\varphi.
\]
Поэтому $E(\mathbf f)=E^*(\mathbf f)=2+\varphi$.
Показатель $1$ не является возможным значением $E^*$ штурмова слова.
\begin{thebibliography}{9}

\bibitem{Diaconis1977}
Diaconis P. The distribution of leading digits and uniform distribution mod~1 //
Ann. Probab. 1977. Vol.~5, No.~1. P.~72--81. DOI: 10.1214/aop/1176995891.
\bibitem{DamanikLenz2002}
Damanik D., Lenz D. The index of Sturmian sequences //
European J. Combin. 2002. Vol.~23. P.~23--29. DOI: 10.1006/eujc.2000.0496.
\bibitem{Peltomaki2015}
Peltomäki J. Characterization of repetitions in Sturmian words: A new proof //
Inform. Process. Lett. 2015. Vol.~115, No.~11. P.~886--891.
DOI: 10.1016/j.ipl.2015.05.011.
\bibitem{AdamczewskiAllouche2007}
Adamczewski B., Allouche J.-P. Reversals and palindromes in continued fractions //
Theoret. Comput. Sci. 2007. Vol.~380. P.~220--237.
DOI: 10.1016/j.tcs.2007.03.017.
\bibitem{TrompShallit1995}
Tromp J., Shallit J. Subword complexity of a generalized Thue--Morse word //
Inform. Process. Lett. 1995. Vol.~54, No.~6. P.~313--316.
DOI: 10.1016/0020-0190(95)00074-M.
\bibitem{Mehraban2006Gulliver}Mehraban E., Gulliver T.A., Hincal E. Blind signatures from the generalized $(t,k)$-Fibonacci $p$-sequences // J. Dyn. Games. 2026. Vol.~13. P.~193--203. DOI: 10.3934/jdg.2025022.
\end{thebibliography}
\end{document}
